\rhead{DS301: Data Visualization \& Storytelling}
\phantomsection 
\section*{Data Visualization \& Storytelling (DS301)}
\label{major:DS_3}

\noindent \small{\hyperref[main:scheme]{$\leftarrow$ Back to Scheme}} \vspace{0.5cm}

\noindent \textbf{Credits:} 3 (2-0-2)

\subsection*{Theory}
\noindent\textbf{Unit 1} \\
\textit{Foundations of Data Visualization:} Data visualization principles and human perception (pre-attentive processing, Gestalt principles), Data types and appropriate encodings (position, length, angle, area, color, shape), Chart taxonomy (categorical, temporal, spatial, hierarchical, network data), Color theory for visualization (perceptual uniformity, color blindness accessibility), Visual encoding principles (Tufte's data-ink ratio, Cleveland/McGill hierarchy).

\vspace{0.3cm}

\noindent\textbf{Unit 2} \\
\textit{Statistical Graphics and Exploratory Analysis:} Exploratory data analysis (EDA) workflow, Univariate distributions (histograms, box plots, violin plots, density plots), Bivariate relationships (scatterplots, heatmaps, correlation matrices), Multivariate visualization techniques (parallel coordinates, scatterplot matrices, dimensionality reduction), Outlier detection and anomaly visualization strategies.

\vspace{0.3cm}

\noindent\textbf{Unit 3} \\
\textit{Geospatial and Network Visualization:} Choropleth maps and spatial aggregation pitfalls, Proportional symbol maps, Flow maps and origin-destination visualization, Spatial autocorrelation and clustering (Moran's I), Network visualization techniques (node-link diagrams, adjacency matrices, arc diagrams), Force-directed layouts, hierarchical edge bundling.

\vspace{0.3cm}

\noindent\textbf{Unit 4} \\
\textit{Interactive and Dynamic Visualizations:} Event handling and brushing/linking, Zooming/panning/filtering interactions, Animated transitions and small multiples, Dashboard design principles (layout, cognitive load, task prioritization), Storytelling techniques (scrollytelling, annotated charts, guided narratives), Level-of-detail management and progressive disclosure.

\vspace{0.3cm}

\noindent\textbf{Unit 5} \\
\textit{Data Storytelling and Presentation:} Narrative structures for data stories (explanatory vs. exploratory analysis), Visual hierarchy and emphasis techniques, Audience analysis and stakeholder communication, Presentation design (slide layouts, chart choice, annotation), Dashboard evaluation frameworks (USE framework, stakeholder interviews), Ethical considerations (visual manipulation, statistical fallacies, misleading scales).

\newpage

\subsection*{Practical}
\noindent\textbf{Lab 1: Univariate and Bivariate Charts} \\
Focuses on creating histograms, box plots, scatterplots, and correlation matrices using proper visual encodings and accessibility considerations.

\vspace{0.2cm}

\noindent\textbf{Lab 2: Multivariate EDA Visualizations} \\
Focuses on parallel coordinates, scatterplot matrices, and dimensionality reduction visualizations for pattern discovery.

\vspace{0.2cm}

\noindent\textbf{Lab 3: Geospatial Data Visualization} \\
Focuses on choropleth maps, proportional symbols, and flow maps using real-world geographic datasets.

\vspace{0.2cm}

\noindent\textbf{Lab 4: Network Graph Visualization} \\
Focuses on node-link diagrams, force-directed layouts, and network metrics visualization.

\vspace{0.2cm}

\noindent\textbf{Lab 5: Interactive Chart Development} \\
Focuses on implementing brushing, filtering, and zooming interactions with D3.js or similar libraries.

\vspace{0.2cm}

\noindent\textbf{Lab 6: Animated Transitions and Storytelling} \\
Focuses on smooth transitions between visualization states and scrollytelling implementations.

\vspace{0.2cm}

\noindent\textbf{Lab 7: Dashboard Design and Layout} \\
Focuses on creating multi-chart dashboards with proper layout, responsive design, and interaction coordination.

\vspace{0.2cm}

\noindent\textbf{Lab 8: Data Story Development} \\
Focuses on complete narrative development from raw data through publication-ready visualizations.

\vspace{0.2cm}

\noindent\textbf{Lab 9: Presentation and Stakeholder Communication} \\
Focuses on presentation design, audience adaptation, and receiving stakeholder feedback.

\vspace{0.2cm}

\noindent\textbf{Lab 10: Visualization Critique and Improvement} \\
Focuses on analyzing existing visualizations, identifying flaws, and redesigning for better clarity and impact.

\vspace{0.2cm}

\newpage
\rhead{Curriculum Schema}
